\section{Results and Discussion}
\label{sec:results-discussion}

Unless stated otherwise, all results use the reference data set of
\cref{subsec:data}: $n = 100$, $\sigma = 0.1$ (so $\sigma^2 = 0.01$ is the
floor of the expected test MSE), an 80/20 split with seed 42, and features
standardized on the training rows. Polynomial degrees run up to $p = 20$.

\subsection{Ordinary least squares}
\label{subsec:results-ols}

\Cref{fig:ols_mse_r2_vs_degree} shows the training and test error of OLS as a
function of the degree. The training MSE decreases monotonically, as it must for
nested least-squares models, from $0.10$ at $p=1$ to $0.0078$ at $p=20$, below
$\sigma^2$: the high-degree fits begin to follow the noise. The test MSE falls
steeply up to $p\approx8$, where the polynomial is first flexible enough to
represent the narrow peak of $f$ at $x = 0$, is smallest at $p=12$ ($0.0086$,
$R^2 = 0.63$), and then rises slowly to $0.0105$ at $p=20$. This is the shape of
fig.~2.11 of \textcite{hastie_2017}, but with a much
flatter right branch than one might expect from Runge's phenomenon. That the
minimum lies slightly below $\sigma^2$ is not a contradiction: with 20 test
points the test MSE has a standard error of roughly
$\sigma^2\sqrt{2/20}\approx0.3\,\sigma^2$, which also makes the location of the
minimum between $p\approx10$ and $p = 14$ essentially arbitrary. The low $R^2$
reflects that most of the variance of $y$ in the test set is noise; the negative
$R^2$ at low degree occurs because the 20 test points happen to contain few
points near the peak, so their variance is small.

\begin{figure}[htbp]
    \centering
    \includegraphics[width=\linewidth]{figures/ols/ols_mse_r2_vs_degree.pdf}
    \caption{OLS training and test MSE (top) and $R^2$ (bottom) vs.\ polynomial
    degree for the reference data set ($n=100$, $\sigma=0.1$, 80/20 split). The
    dotted line is $\sigma^2$; the dashed line marks the minimum test MSE
    ($p=12$).}
    \label{fig:ols_mse_r2_vs_degree}
\end{figure}

The coefficients, shown in \cref{fig:ols_coefficients_vs_degree}, tell the other half of
the story. Up to $p\approx6$ they are of order one, but from there on they grow
by roughly a factor of three every two degrees, to $\max_j|\hat\theta_j| = 3.3\times
10^3$ at $p=20$, with alternating signs between neighboring powers. The fitted
curve is a near-cancellation of large, nearly collinear terms: the condition
number of the standardized design matrix grows from $25$ at $p=5$ to $1.7\times
10^{3}$ at $p=10$ and $1.6\times10^7$ at $p=20$. Such coefficients are
determined by the particular noise realization, which is the variance that the
bias--variance analysis below quantifies.

\begin{figure}[htbp]
    \centering
    \includegraphics[width=\linewidth]{figures/ols/ols_coefficients_vs_degree.pdf}
    \caption{OLS coefficients $\hat\theta_j$ of the standardized feature $x^j$
    (column) for each degree $p$ (row), on a symmetric logarithmic color scale.
    The intercept, equal to the mean training target, is not shown.}
    \label{fig:ols_coefficients_vs_degree}
\end{figure}

With more data, the minimum moves to higher degree and the high-degree
deterioration disappears (\cref{fig:ols_test_mse_by_n}): for $n = 400$ the test
MSE is flat at $\approx\sigma^2$ from $p\approx12$ to $20$, whereas for $n=50$
(with only 10 test points) the curve is erratic beyond $p = 12$ and reaches $3.2$
at $p=20$. The erratic curves for $n=50$ and $n=200$ are dominated by single
test points near the ends of the interval, a mechanism we return to in
\cref{subsec:results-selection}.

\begin{figure}[htbp]
    \centering
    \includegraphics[width=\linewidth]{figures/ols/ols_test_mse_by_n.pdf}
    \caption{OLS test MSE vs.\ degree for $n \in \{50, 100, 200, 400\}$ points
    ($\sigma=0.1$, 80/20 split, seed 42 for every $n$).}
    \label{fig:ols_test_mse_by_n}
\end{figure}

\subsection{Ridge regression}
\label{subsec:results-ridge}

\Cref{fig:ridge_shrinkage} shows the mechanism of Ridge regression at $p=15$.
The squared singular values of the standardized design matrix, $s_i^2/n$, span
ten orders of magnitude, from $9.8$ to $3.8\times10^{-10}$. According to
\cref{eq:ridge_scaling}, a penalty $\lambda$ suppresses the modes with
$s_i^2/n \lesssim \lambda$ and leaves the others untouched, so increasing
$\lambda$ by a decade removes roughly one more mode. $\lambda = 10^{-10}$ is
essentially OLS; $\lambda = 1$ keeps only the first three modes. The ridge
trace (\cref{fig:ridge_trace}) shows the corresponding coefficients: the large,
alternating OLS coefficients collapse between $\lambda\approx10^{-9}$ and
$10^{-6}$, as the smallest modes are removed.

\begin{figure}[htbp]
    \centering
    \includegraphics[width=\linewidth]{figures/ridge/ridge_shrinkage_degree15.pdf}
    \caption{Ridge shrinkage factors $d_i = s_i^2/(s_i^2+n\lambda)$
    (\cref{eq:ridge_scaling}) of the SVD modes of the standardized design matrix at
    $p=15$ (training set of the reference data), for $\lambda$ from $10^{-10}$
    to $1$. Dotted, right axis: $s_i^2/n$.}
    \label{fig:ridge_shrinkage}
\end{figure}

\begin{figure}[htbp]
    \centering
    \includegraphics[width=\linewidth]{figures/ridge/ridge_trace_degree15.pdf}
    \caption{Ridge trace at $p=15$: coefficients $\hat\theta_j(\lambda)$ of the
    standardized features, colored by the power $j$.}
    \label{fig:ridge_trace}
\end{figure}

\Cref{fig:ridge_test_mse} repeats the analysis of \cref{fig:ols_mse_r2_vs_degree}
for Ridge. A small penalty ($\lambda = 10^{-8}$) reproduces the OLS curve up to
$p\approx14$ and removes the slight rise beyond it ($0.0096$ instead of $0.0105$
at $p = 20$); $\lambda = 10^{-6}$ is similar. Larger penalties make the curve
flat in $p$, because the extra degrees only add modes that the penalty
suppresses, but at a higher level: $0.013$ for $\lambda = 10^{-4}$ and
$0.023$ for $\lambda = 10^{-2}$. Here the penalty removes modes that are needed
to represent the peak of $f$, and the bias dominates. On this data set Ridge
therefore makes the choice of degree less critical, but it does not lower the
best achievable test error; we test this more carefully in
\cref{subsec:results-selection}.

\begin{figure}[htbp]
    \centering
    \includegraphics[width=\linewidth]{figures/ridge/ridge_test_mse_by_lambda.pdf}
    \caption{Ridge test MSE vs.\ degree for four values of $\lambda$, compared
    with OLS (reference data set, single split). Dotted: $\sigma^2$.}
    \label{fig:ridge_test_mse}
\end{figure}

\subsection{Bias--variance trade-off}
\label{subsec:results-bias-variance}

\Cref{fig:bias_variance_tradeoff_bootstrap} shows the bootstrap decomposition
(\cref{eq:bias-variance-decomposition}) for OLS with $B = 100$ resamples of the
training set, evaluated on the fixed test set. At low degree the error is
dominated by the bias: at $p = 2$ the squared bias against $f$ is $0.032$ and the
variance $0.002$. Increasing the degree trades bias for variance, and the
bootstrap test error is smallest at $p = 6$ ($0.021$), where
$\mathrm{Bias}^2 = 0.009$ and $\mathrm{Var} = 0.006$ are of the same size.
Beyond it the variance explodes, from $0.006$ at $p=6$ to $0.57$ at $p=10$ and
$15$ at $p = 15$: a factor of $2.6\times10^3$ between $p=6$ and $p=15$. The two
bias estimates differ by approximately $\sigma^2$, as
\cref{eq:measured-bias} predicts (e.g.\ $0.015$ against $\bm{y}$ and $0.009$
against $\bm{f}$ at $p=6$; the difference fluctuates around $\sigma^2 = 0.01$
since only 20 test points enter the average).

\begin{figure}[htbp]
    \centering
    \includegraphics[width=\linewidth]{figures/bias_variance/bias_variance_tradeoff_bootstrap.pdf}
    \caption{Bootstrap bias--variance decomposition of the OLS test error
    ($B=100$, reference data set). The squared bias is measured both against the
    noisy test targets $\bm{y}$ (\cref{eq:measured-bias}, absorbs $\sigma^2$) and
    against the true $\bm{f}$ (\cref{eq:bias-term}). Dashed: minimum of the
    bootstrap test MSE ($p = 6$). Values above $3$ are cut off.}
    \label{fig:bias_variance_tradeoff_bootstrap}
\end{figure}

Two features of the figure deserve comment. First, the bootstrap error is
minimized at $p=6$, whereas the single test split of
\cref{fig:ols_mse_r2_vs_degree} favors $p = 12$; both use the same training and
test points. The difference arises because each bootstrap fit sees only about
$63\%$ of the distinct training points (\cref{subsec:resampling}) and is
therefore much less stable at high degree than the fit on all 80 points. The
bootstrap thus estimates the variance of a model trained on fewer, and
duplicated, points. Second, the bias is not monotone in $p$ beyond $p\approx8$:
the average over bootstrap fits is itself contaminated by a few fits that go
wild near the ends of the interval, so $\mathbb{E}[\tilde{\bm y}]$ is not a
smooth polynomial fit either. We examine the origin of these wild fits in
\cref{subsec:results-selection}.

The decomposition responds to $n$ and $\sigma$ as expected
(\cref{fig:bias_variance_by_n,fig:bias_variance_by_noise} in
\cref{app:additional-results}). With more data the variance at a given degree
drops while the bias hardly changes, so the minimum moves from $p = 4$ at $n=50$
to $6$, $10$ and $12$ at $n=100$, $200$ and $400$. More noise raises the floor
and the variance, and the optimal degree falls from $6$ ($\sigma \leq 0.2$) to
$4$ ($\sigma = 0.4$).

\subsection{Resampling and model selection}
\label{subsec:results-selection}

\subsubsection{Which degree do the estimators select?}
\Cref{fig:ols_estimators}(a) compares four estimates of the OLS test error on
the reference data set, all computed from the 80 training points: the bootstrap
(on the held-out 20 points), 5- and 10-fold CV, and the single split. Since our
data are synthetic, we can also compute the \emph{true} expected test error of
the model fitted on all 80 points, $\mathbb{E}_x[(\hat f(x)-f(x))^2] + \sigma^2$
for $x\sim\mathcal{U}(-1,1)$, on a dense grid. The true error is smallest at
$p=14$ ($0.0133$) and stays below $0.026$ up to $p=20$. The four estimates agree
up to $p\approx7$, but then separate: the bootstrap selects $p=6$, both CV
variants $p = 8$ and the single split $p=12$, and the bootstrap and CV
estimates grow by several orders of magnitude at high degree, where the true
error does not.

\begin{figure*}[htbp]
    \centering
    \includegraphics[width=\textwidth]{figures/cross_validation/ols_estimators_vs_degree.pdf}
    \caption{Estimated OLS test error vs.\ degree. (a) Reference data set:
    bootstrap ($B = 100$), 5- and 10-fold CV and a single 80/20 split, all
    computed from the same 80 training points, and the true error of the fit on
    those 80 points (from the known $f$). Open markers: minima. (b) Degree
    selected by each estimator over 31 independent data sets.}
    \label{fig:ols_estimators}
\end{figure*}

\Cref{fig:cv_extrapolation} shows where the large estimates come from. In the
5-fold split that dominates the CV estimate at $p = 15$, the training folds
happen to contain no point beyond $x = 0.854$, while the held-out fold contains
the largest training input, $x = 0.951$. The degree-15 polynomial fitted on the
other folds must extrapolate to that point and predicts $-34.8$ there, while
the fit on all 80 points is well behaved. This single point contributes $87\%$
of the squared error of that fold and makes the fold MSE $88$. The mechanism
is the statistical counterpart of Runge's phenomenon: a high-degree polynomial
is well constrained only where it has data, and it diverges quickly just
outside that range. It affects the bootstrap even more: for the reference data,
$36\%$ of the bootstrap samples leave at least one test point outside their
$x$-range. The resampling estimates thus include extrapolation errors that the
final model, trained on all points, does not make. Leave-out-based estimators
are therefore pessimistic for high-degree polynomials, and increasingly so the
fewer points each fit uses (bootstrap $>$ 5-fold $>$ 10-fold).

\begin{figure}[htbp]
    \centering
    \includegraphics[width=\linewidth]{figures/cross_validation/cv_extrapolation_fold.pdf}
    \caption{The 5-fold CV split that dominates the CV estimate at $p=15$
    (reference data). Orange: fit on the four training folds (gray points),
    which must extrapolate beyond their range (shaded) to the held-out point at
    $x = 0.951$. Blue: fit on all 80 training points.}
    \label{fig:cv_extrapolation}
\end{figure}

\Cref{fig:ols_estimators}(b) shows that this is systematic. Over 31 data sets,
the median selected degree is $8$ for the bootstrap, $10$ for 5-fold and $12$ for
10-fold CV, and $13$ for the single split, while the degree with the smallest
true error has median $14$. The selection matters less than the spread
suggests, because the true error is flat over a wide range of degrees: the mean
true error of the model at the selected degree is $0.0179\pm0.0008$ for the
bootstrap, $0.024\pm0.006$ for 5-fold CV, $0.028\pm0.008$ for 10-fold CV and
$0.022\pm0.006$ for the single split (mean $\pm$ standard error over the 31 data
sets), against $0.0131\pm0.0003$ at the best possible degree. The bootstrap,
although it selects too low a degree, is the most \emph{reliable} selector: CV
and the single split occasionally select a degree whose fit extrapolates badly
at the ends of the interval, which is what produces their large standard
errors. The two resampling methods therefore do not point to the same model
complexity, and neither points to the best one, for a reason specific to
polynomials: their error estimates are dominated by extrapolation.

\subsubsection{Ridge and Lasso over degree and penalty}
\Cref{fig:cv_heatmaps} shows the CV error of Ridge and Lasso over the grid of
degrees and penalties on the reference training set. The Ridge surfaces have a
valley running diagonally: higher degrees need larger penalties. In their lower
right corner, small $\lambda$ at high degree, the CV error explodes as for OLS,
because the penalty is too weak to prevent the extrapolation failures of
\cref{fig:cv_extrapolation}; $\lambda \gtrsim 10^{-7}$ suppresses them. 5-fold
CV selects Ridge with $p = 11$, $\lambda = 10^{-6}$ (CV MSE $0.0121$) and
10-fold CV $p = 12$, $\lambda = 10^{-5.5}$ ($0.0104$), whereas OLS with 5-fold
CV selects $p = 8$ with a CV MSE of $0.0285$. The Ridge minimum is less than
half the OLS minimum, but on the same five folds the difference is $-0.016 \pm
0.013$ (mean $\pm$ standard error over folds), not significant: Ridge mainly
avoids the one fold with the extrapolation failure that inflates the OLS
estimate. The Lasso surface looks different. It hardly depends on the degree
beyond $p\approx8$ and has its minimum at a low degree and a large penalty, $p
= 6$ and $\lambda = 10^{-2.5}$ (CV MSE $0.0222$). As \cref{subsec:results-lasso}
shows, this reflects our solver more than the Lasso model: with 5000 subgradient
steps, the high-degree fits are far from the Lasso solution.

\begin{figure*}[htbp]
    \centering
    \includegraphics[width=\textwidth]{figures/cross_validation/cv_heatmaps.pdf}
    \caption{CV test MSE over polynomial degree and penalty on the reference
    training set: Ridge with (a) 5-fold and (b) 10-fold CV, and (c) Lasso (our
    subgradient solver, Adam, $\gamma = 0.05$, 5000 iterations) with 5-fold CV.
    Stars: minima.
    Colors are clipped at $1$.}
    \label{fig:cv_heatmaps}
\end{figure*}


\subsection{Gradient descent}
\label{subsec:results-gd}

Linear regression is an ideal test case for gradient descent: the closed-form
solutions are exact benchmarks, the gradient is known analytically, and the
convexity of the cost guarantees convergence for a small enough learning rate.
We compare the optimizers at $p = 5$, where plain GD converges in a reasonable
number of iterations, and use $p = 10$ and $15$ to study ill-conditioning. All
runs start from $\bm{\theta} = \bm{0}$.

\Cref{fig:gd_convergence} shows plain GD with $\gamma = 1/L$ at $p=5$, where
$\kappa(\bm{H}) = 604$ for OLS and $510$ for Ridge with $\lambda = 10^{-3}$.
The relative cost gap decreases linearly on the logarithmic scale, as the
eigen-analysis of \cref{subsec:gradient-descent} predicts: after 6000 iterations
it is $10^{-10}$ for OLS and $10^{-12}$ for Ridge, whose better conditioning
gives the steeper slope. The runs with analytical and automatically
differentiated gradients are indistinguishable, as expected from their agreement
to $5\times10^{-16}$ (\cref{subsec:verification}).

\begin{figure}[htbp]
    \centering
    \includegraphics[width=\linewidth]{figures/gradient_descent/gd_convergence.pdf}
    \caption{Plain gradient descent with $\gamma = 1/L$ at $p=5$ (reference
    training set): relative gap between the cost and the cost of the closed-form
    solution vs.\ iteration, for OLS and Ridge ($\lambda=10^{-3}$), with
    analytical (thick) and JAX (dashed) gradients.}
    \label{fig:gd_convergence}
\end{figure}

\Cref{fig:gd_stability} tests the stability bound \cref{eq:gd_stability}. After
2000 steps the cost gap decreases steadily as $\gamma$ approaches
$2/\lambda_{\max}(\bm{H})$ and jumps to divergence immediately beyond it, for
both OLS and Ridge. The best fixed learning rate lies just below the bound:
the optimal constant step for a quadratic is $2/(\lambda_{\max} +
\lambda_{\min})$, which is within $0.2\%$ of $2/\lambda_{\max}$ when
$\kappa \approx 600$.

\begin{figure}[htbp]
    \centering
    \includegraphics[width=\linewidth]{figures/gradient_descent/gd_stability_sweep.pdf}
    \caption{Relative cost gap after 2000 plain-GD steps vs.\ the learning rate
    in units of the bound $2/\lambda_{\max}(\bm{H})$ ($p = 5$). Diverged runs are
    shown at $10^{12}$.}
    \label{fig:gd_stability}
\end{figure}

\Cref{tab:optimizers} compares the five update rules by the number of
iterations needed to reach $\lVert\bm{\theta}-\hat{\bm{\theta}}\rVert_2 /
\lVert\hat{\bm{\theta}}\rVert_2 < 10^{-3}$, with the learning rate tuned for
each optimizer and problem on a logarithmic grid of 25 values between $10^{-5}$
and $10$ (four per decade), and a cap of $10^5$ iterations.
\Cref{fig:optimizer_lr_sensitivity} shows the full dependence on the learning
rate for two of the problems. Three observations stand out.

First, momentum and Adam are an order of magnitude faster than the other
update rules. At $p=5$, momentum needs $126$ iterations and Adam $337$, against
$3166$ for plain GD, a factor of $25$ and $9.4$; AdaGrad and RMSprop are no
faster than plain GD. For the Ridge problems at $p = 10$ and $15$ the two are
equally fast within the resolution of the learning-rate grid ($1342$ vs.\
$1422$ and $1913$ vs.\ $1379$ iterations), about $18$ times faster than plain
GD. Momentum also converges for learning rates beyond the plain-GD bound
(\cref{fig:optimizer_lr_sensitivity}a); this agrees with the stability bound of
the heavy-ball iteration for a quadratic, $\gamma < 2(1+\beta)/\lambda_{\max}
= 0.58$ at $p=5$.

Second, Adam is by far the least sensitive to the learning rate: it reaches the
tolerance for all 25 learning rates at $p=5$ and for 23 of 25 at $p = 10$ and
$15$, and for $\gamma \gtrsim 0.03$ its iteration count is essentially
independent of $\gamma$. Plain GD converges for only 3--7 of the 25 values, and
momentum for 8--13; both need $\gamma$ within a factor of a few of the bound
\cref{eq:gd_stability}. AdaGrad and Adam are insensitive to large $\gamma$
because they normalize the step by the gradient magnitude. RMSprop, which does
the same with a short memory, only converges for $\gamma \leq 6\times10^{-4}$:
at larger $\gamma$ its normalized steps keep it oscillating around the minimum.

Third, the conditioning decides everything. For OLS at $p = 10$ and $15$, with
$\kappa(\bm{H}) = 2.8\times10^6$ and $2.6\times10^{10}$, no optimizer reaches
the tolerance within $10^5$ iterations; the best relative error at $p=10$ is
$0.28$ (Adam), and at $p = 15$ none gets below $0.99$, barely closer to
$\hat{\bm{\theta}}$ than the starting point $\bm{\theta}=\bm{0}$ (error $1$). With $\gamma = 1/L$, plain GD
would need about $\kappa\ln 10^3 \approx 2\times10^7$ iterations at $p=10$. A
penalty of only $\lambda = 10^{-3}$, which barely changes the fit, reduces
$\kappa$ to $\sim10^4$ (\cref{subsec:gradient-descent}), and all optimizers
converge again. Where they converge, the iteration counts of plain GD scale
with $\kappa$ as predicted: $3166 \approx 5\kappa$ at $p=5$ and $26244 \approx 4
\kappa$ for Ridge at $p = 10$. The OLS parameters at high degree are, however,
ill-determined anyway (\cref{fig:ols_coefficients_vs_degree}): a relative
parameter error of order one does not preclude a good prediction.


\begin{table*}[htbp]
    \centering
    \caption{Iterations to reach the closed-form OLS or Ridge ($\lambda =
    10^{-3}$) solution to a relative parameter error of $10^{-3}$, from
    $\bm{\theta} = \bm{0}$, with the learning rate tuned per optimizer and
    problem (reference training set, standardized features, degree $p$).
    $\kappa(\bm{H})$ is the condition number of the Hessian. In parentheses: the
    number of the 25 learning rates for which the tolerance was reached within
    $10^5$ iterations. ``--'': not reached for any learning rate; in brackets the
    smallest relative error reached. Bold: fastest.}
    \label{tab:optimizers}
    \input{figures/tables/optimizers.tex}
\end{table*}

\begin{figure*}[htbp]
    \centering
    \includegraphics[width=\textwidth]{figures/gradient_descent/optimizer_lr_sensitivity.pdf}
    \caption{Iterations to reach the closed-form solution to a relative
    parameter error of $10^{-3}$ vs.\ the learning rate, for each optimizer.
    Missing points: tolerance not reached within $10^5$ iterations, or
    divergence. Dotted: the plain-GD stability bound $2/\lambda_{\max}(\bm{H})$.}
    \label{fig:optimizer_lr_sensitivity}
\end{figure*}


\subsection{Lasso}
\label{subsec:results-lasso}

We solve Lasso with our gradient-descent code, using the subgradient
\cref{eq:lasso_subgradient} with each of the update rules of
\cref{subsec:results-gd}, and compare with scikit-learn's coordinate-descent
\texttt{Lasso}, whose solution defines the optimal cost $C^*$.
\Cref{fig:lasso_convergence} shows the convergence at two configurations of the
reference data: the one selected by 5-fold CV ($p = 6$, $\lambda =
10^{-2.5}$), and a high-degree one ($p = 15$, $\lambda = 10^{-4}$).

\begin{figure*}[htbp]
    \centering
    \includegraphics[width=\textwidth]{figures/lasso/lasso_convergence.pdf}
    \caption{Relative gap between the Lasso cost of subgradient descent and the
    cost $C^*$ of scikit-learn's solution vs.\ iteration, from $\bm{\theta} =
    \bm{0}$, for each update rule (reference training set). (a) The
    configuration selected by 5-fold CV, $p=6$, $\lambda=10^{-2.5}$. (b) $p=15$,
    $\lambda=10^{-4}$. Plain GD and momentum use $\gamma = 1/L$.}
    \label{fig:lasso_convergence}
\end{figure*}

At $p=6$, where the smooth part of the cost has $\kappa = 2.3\times10^3$, all
update rules approach $C^*$, but none converges in the usual sense. Each levels
off at a gap set by its step size, as expected for the subgradient method with
a fixed step (\cref{subsec:gradient-descent}): momentum at $8\times10^{-6}$
after a few thousand iterations, RMSprop at $4\times10^{-4}$, and Adam, whose
normalized step does not shrink near the minimum, oscillates between
$10^{-4}$ and $10^{-2}$ and is at $2.4\times10^{-3}$ after $5\times10^{4}$
iterations, worse than after 5000 ($8.7\times10^{-4}$). Plain GD and AdaGrad,
whose effective steps are smallest near the end, reach $10^{-6}$. The
coefficients from our default setting (Adam, 5000 iterations) differ from
scikit-learn's by at most $4\times10^{-3}$, but scikit-learn sets the
coefficient of $x^5$ exactly to zero, while ours is $-3.8\times10^{-3}$: our
solver produced no exact zeros in any of the fits we examined.

The analytical subgradient uses $\operatorname{sign}(0)=0$; the gradient from
JAX of $|\theta_j|$ at exactly $\theta_j = 0$ is $1$ (\cref{subsec:verification}).
Both are valid subgradients, since $\partial|\theta_j| = [-1,1]$ at zero, but
the choice rarely matters in practice: the iterates are almost never exactly
zero, which is the reason for the oscillation in the first place.

At $p = 15$, with $\kappa = 2.6\times10^{10}$, the picture is much worse
(\cref{fig:lasso_convergence}b). After $5\times10^4$ iterations Adam is at a
gap of $1.2\times10^{-2}$, momentum at $2.5\times10^{-2}$, and the others above
$0.1$; scikit-learn's coordinate descent itself needs $4\times10^5$ passes.
\Cref{fig:lasso_coefficients} compares the coefficients: scikit-learn sets five
of the 15 to zero, while ours after 5000 iterations differ by up to $4.8$ and
are as large as $1.7$ where scikit-learn's are zero; even after $5\times10^4$
iterations they differ by up to $0.5$. Our Lasso fits at high degree are
therefore early-stopped gradient descent rather than Lasso solutions. This
explains the Lasso CV surface of \cref{fig:cv_heatmaps}c: at high degree the
fits are poor regardless of $\lambda$, so CV prefers a low degree, where the
solver converges.

\begin{figure}[htbp]
    \centering
    \includegraphics[width=\linewidth]{figures/lasso/lasso_coefficients.pdf}
    \caption{Lasso coefficients of the standardized features at $p=15$,
    $\lambda=10^{-4}$ (reference training set): scikit-learn (circles; squares
    mark exact zeros) and our subgradient descent with Adam after 5000 and
    $5\times10^4$ iterations.}
    \label{fig:lasso_coefficients}
\end{figure}

\subsection{Stochastic gradient descent}
\label{subsec:results-sgd}

\Cref{fig:sgd_batch} and \cref{tab:sgd} compare mini-batch SGD with full-batch
GD for OLS at $p = 5$, at equal computational cost measured in floating-point
operations of the gradient evaluations: one epoch of SGD costs as much as one
full-batch step. Plain SGD uses $\gamma_0 = 1/L$ and Adam $\gamma_0=0.03$ (the
best fixed rate of \cref{tab:optimizers}), both with the decay $\gamma_t =
\gamma_0/(1+10^{-3}t)$.

With plain updates, small mini-batches make faster progress per FLOP early on,
because they make $80/M$ updates per epoch instead of one: after 200 epochs the
relative cost gap is $2.7\times10^{-3}$ with $M = 4$ and $2.3\times10^{-2}$ with
the full batch. Later, the gradient noise sets a floor: all batch sizes end
between $9\times10^{-5}$ and $6\times10^{-4}$ after 2000 epochs, fluctuating
from epoch to epoch, while full-batch GD reaches $5.1\times10^{-5}$ monotonically.
With Adam the picture is reversed: full-batch Adam reaches the closed-form
cost to machine precision within about $10^6$ FLOPs, whereas mini-batch Adam
stalls at $10^{-5}$--$10^{-3}$. Full-batch Adam also shows the behavior noted
in \cref{subsec:gradient-descent}: after converging, its fixed step size
produces a brief excursion (the spike in \cref{fig:sgd_batch}b) before it
settles again.

\begin{figure*}[htbp]
    \centering
    \includegraphics[width=\textwidth]{figures/sgd/sgd_batch_size_sweep.pdf}
    \caption{Relative cost gap (OLS, $p=5$) vs.\ cumulative FLOPs of the gradient
    evaluations for mini-batch sizes $M \in \{4, 16, 32\}$ and the full batch
    ($M = 80$), over 2000 epochs, for (a) plain SGD with $\gamma_0 = 1/L$ and (b)
    Adam with $\gamma_0 = 0.03$; SGD with the decay $\gamma_t=\gamma_0/(1+10^{-3}t)$.
    Values below $10^{-16}$ are shown at $10^{-16}$.}
    \label{fig:sgd_batch}
\end{figure*}

\begin{table}[htbp]
    \centering
    \caption{Full-batch GD vs.\ mini-batch SGD ($M = 16$) after 2000 epochs at
    $p=5$: number of updates, FLOPs of the gradient evaluations, final relative
    cost gap, and median wall-clock time of three runs.}
    \label{tab:sgd}
    {\footnotesize\setlength{\tabcolsep}{3pt}
    \input{figures/tables/sgd.tex}}
\end{table}

\Cref{tab:sgd} also shows the price of SGD in practice. With equal FLOPs, SGD
takes three to four times longer in wall-clock time, because each of its five
times more numerous updates carries a fixed Python overhead that is large
compared with the arithmetic for 16 rows. For a problem with 80 training points
and 6 parameters, SGD thus has no advantage: its benefit, cheap updates when
$n$ is large and the gradient is expensive, does not apply, and its noise
prevents it from reaching the accuracy of full-batch methods without a carefully
tuned decay of the learning rate.


\subsection{Final comparison of OLS, Ridge and Lasso}
\label{subsec:results-final}

The comparison on the reference data rests on one set of folds and one
test set of 20 points. To compare OLS, Ridge and Lasso reliably, we repeated
the whole selection on 31 independent data sets: 5-fold CV on the 80 training
points over the grids of \cref{fig:cv_heatmaps}, a refit of the selected
configuration on all 80 points, and evaluation on the 20 held-out points and,
using the known $f$, of the true error. \Cref{fig:final_comparison} and
\cref{tab:final} summarize the results.

\begin{figure*}[htbp]
    \centering
    \includegraphics[width=\textwidth]{figures/cross_validation/final_comparison.pdf}
    \caption{Error of the model selected by 5-fold CV for each method on 31
    independent data sets ($n=100$, $\sigma = 0.1$): (a) on the 20 held-out
    test points, (b) true error from the known $f$. Gray lines connect the
    same data set; black bars are medians; dotted: $\sigma^2$.}
    \label{fig:final_comparison}
\end{figure*}

\begin{table}[htbp]
    \centering
    \caption{Model selection over 31 data sets: median CV, test and true MSE of
    the CV-selected model, median selected degree, and the number of data sets
    in which the method beats OLS.}
    \label{tab:final}
    {\footnotesize\setlength{\tabcolsep}{4pt}
    \begin{tabular}{@{}lccccc@{}}
        \toprule
        & CV MSE & test MSE & true MSE & $\hat p$ & better than OLS \\
        & \multicolumn{3}{c}{(median)} & (median) & CV/test/true \\
        \midrule
        OLS & $0.0139$ & $0.0134$ & $0.0146$ & 10 & -- \\
        Ridge & $0.0129$ & $0.0136$ & $0.0137$ & 12 & 25/15/19 \\
        Lasso & $0.0150$ & $0.0134$ & $0.0156$ & 8 & 11/12/10 \\
        \bottomrule
    \end{tabular}}
\end{table}

Judged by the CV error itself, Ridge beats OLS in 25 of the 31 data sets, by a
median of $5\%$. This is the number a naive comparison would report, and it is
biased: the CV minimum of Ridge is taken over a two-dimensional grid, that of
OLS over a one-dimensional one, and the minimum of noisy estimates over more
candidates is lower even when no candidate is better. On held-out data the
advantage disappears: Ridge beats OLS on the test set in 15 of 31 data sets,
with a median relative difference of $0.0\%$, and in the true error in 19 of
31, with a median of $-0.2\%$. The penalty buys stability rather than
accuracy: it lets CV select a higher degree (median 12 instead of 10), with a
smaller bias, without an increase in variance. The means are dominated by a few
data sets in which the selected model extrapolates badly: Ridge with $p=15$ and
$\lambda = 10^{-10}$ reaches a true MSE of $18$ on one data set, where all of
the error comes from $x$ beyond the largest training input.

Lasso, fitted with our subgradient solver, is worse than OLS by every measure:
it beats OLS in only 11 (CV), 12 (test) and 10 (true error) of the 31 data sets,
with a median relative increase of its true error of $13\%$. CV selects a low degree for it
(median 8), for the reason found in \cref{subsec:results-lasso}: at high degree
the solver does not reach the Lasso solution within 5000 iterations, so these
fits are poor and CV avoids them. None of the selected Lasso models has an exact
zero coefficient. This comparison therefore measures our optimizer as much as
the Lasso model; a solver that converges at high degree could change the
ranking.

Which term of the bias--variance decomposition does each method attack? OLS
controls the variance only through the degree. Ridge reduces the variance
directly by shrinking the small-singular-value directions
(\cref{fig:ridge_shrinkage}), at the cost of bias along them, which allows a
more flexible, less biased model for the same variance. Lasso aims at the same
term by setting coefficients exactly to zero, but our fits never reach that
regime; what limits them in practice is the incomplete optimization, which acts
like early stopping and adds bias. Our numerical evidence supports the reading
for Ridge qualitatively (higher selected degrees, flat error curves in
\cref{fig:ridge_test_mse}), but it also shows its limits for this problem. The
selected Ridge penalties are tiny, $\lambda \in [10^{-12}, 3\times10^{-5}]$
(5 of 31 at the lower end of the grid, i.e.\ effectively OLS), because with $n
= 100$ and $\sigma=0.1$ the remaining error of a well-chosen OLS fit is
dominated not by variance in the interior of the interval but by the
extrapolation near its ends, which a global penalty barely touches. For Runge's
function with a random design, no method is thus significantly better than OLS
on the held-out data we would have in practice, and Ridge is better by less than
$1\%$ in the true error.
